
\documentclass{article} % For LaTeX2e
\usepackage{iclr2026_conference,times}

% Optional math commands from https://github.com/goodfeli/dlbook_notation.
\input{math_commands.tex}

\usepackage{hyperref}
\usepackage{url}
\usepackage{booktabs}       % Professional-quality tables
\usepackage{amsfonts}       % Blackboard math symbols
\usepackage{nicefrac}       % Compact symbols for 1/2, etc.
\usepackage{microtype}      % Microtypography
\usepackage{graphicx}
\usepackage{xcolor}

% ---------------------------------------------------------------
% PAPER METADATA
% ---------------------------------------------------------------
\title{[Paper Title]}

% For blind review: do NOT include author information.
% Uncomment \iclrfinalcopy and fill in authors for camera-ready.

% \author{Author One$^{1}$ \quad Author Two$^{1}$ \quad Author Three$^{2}$ \\
% $^{1}$Institution One \quad $^{2}$Institution Two \\
% \texttt{\{author1, author2\}@institution1.edu} \quad
% \texttt{author3@institution2.edu}
% }

% Uncomment for camera-ready version ONLY:
%\iclrfinalcopy

\newcommand{\method}{[MethodName]}   % replace with your method name

% ---------------------------------------------------------------
\begin{document}

\maketitle

% ---------------------------------------------------------------
\begin{abstract}
% 5-sentence formula (Farquhar / DeepMind):
% 1. What you achieved: "We introduce / prove / demonstrate..."
% 2. Why this is hard and important
% 3. How you do it (with specialist keywords for discoverability)
% 4. What evidence you have
% 5. Your most remarkable number / result
%
% DELETE the line above before submission.
[Abstract placeholder. Replace with 5-sentence abstract following the formula above.]
\end{abstract}

% ---------------------------------------------------------------
\section{Introduction}

% Goal: ≤ 1.5 pages. Must end with a contribution bullet list.
%
% Structure:
%   - Problem statement (what problem do we solve and why it matters)
%   - Gap in existing work (what is missing)
%   - Our approach (brief, 2–3 sentences)
%   - Contribution list (2–4 bullets, 1–2 lines each)
%   - Paper organization (optional, 1 sentence per section)

[Introduction placeholder.]

\textbf{Contributions.} Our main contributions are:
\begin{itemize}
  \item We propose \method{}, which [contribution 1].
  \item We demonstrate [contribution 2 / key result].
  \item We provide [contribution 3, e.g., theoretical analysis / open-source code].
\end{itemize}

% ---------------------------------------------------------------
\section{Related Work}

% Organize methodologically, NOT paper-by-paper.
% Good: "One line of work assumes X [refs]; we instead assume Y because..."
% Bad: "A et al. do X. B et al. do Y."
% Cite generously — reviewers likely authored relevant work.

[Related work placeholder. Organize by research theme, not chronologically.]

% ---------------------------------------------------------------
\section{Method}

% Enable reimplementation:
%   - Conceptual outline or pseudocode
%   - All hyperparameters listed explicitly
%   - Architectural details sufficient for reproduction

\subsection{Problem Setup}

[Formal problem definition. State all assumptions before any theorem or claim.]

\subsection{\method{}}

[Core method description. Use pseudocode or algorithm environment if helpful.]

% ---------------------------------------------------------------
\section{Experiments}

% For each experiment, state:
%   1. What claim it supports
%   2. How it connects to the main contribution
%   3. Experimental setting (details may go to Appendix)
%   4. What to observe: "the blue line shows X, which demonstrates Y"
%
% Requirements: error bars, \#seeds, compute (GPU type + hours), baselines.

\subsection{Experimental Setup}

[Datasets, baselines, evaluation metrics, compute resources.]

\subsection{Main Results}

[Main comparison table. Use booktabs. Bold best values. Include ↑/↓ for metrics.]

\begin{table}[t]
  \caption{[Table caption. Lower case except first word and proper nouns.]}
  \label{tab:main}
  \centering
  \begin{tabular}{lccc}
    \toprule
    Method & Metric A ↑ & Metric B ↑ & Metric C ↓ \\
    \midrule
    Baseline 1   & 00.0 & 00.0 & 00.0 \\
    Baseline 2   & 00.0 & 00.0 & 00.0 \\
    \midrule
    \textbf{\method{} (Ours)} & \textbf{00.0} & \textbf{00.0} & \textbf{00.0} \\
    \bottomrule
  \end{tabular}
\end{table}

\subsection{Ablation Studies}

[Ablation table. Each row removes / modifies one component.]

% ---------------------------------------------------------------
\section{Discussion}

[Interpretation of results. Connection to the narrative / main claims.]

% ---------------------------------------------------------------
\section{Limitations}

% REQUIRED by ICLR. Counter-intuitively, honesty helps:
%   - Reviewers are instructed NOT to penalize honest limitations
%   - Pre-empt criticisms by identifying weaknesses first
%   - Explain why limitations do not undermine the core claims

[Discuss at least 2–3 genuine limitations of the work.]

% ---------------------------------------------------------------
\section{Conclusion}

[One paragraph. State the core contribution, main finding, and one forward-looking sentence.
Do NOT introduce new claims here.]

% ---------------------------------------------------------------
% LLM USAGE DISCLOSURE (required by ICLR 2026)
% ---------------------------------------------------------------
\subsubsection*{LLM Usage Disclosure}

[State which LLMs (if any) were used, and for what purpose
(e.g., grammar checking, code generation). Required by ICLR 2026.]

% ---------------------------------------------------------------
\subsubsection*{Acknowledgments}

Use unnumbered third-level headings for acknowledgments.
All acknowledgments, including funding agencies, go here.

% ---------------------------------------------------------------
\bibliography{references}
\bibliographystyle{iclr2026_conference}

% ---------------------------------------------------------------
\appendix

\section{Implementation Details}

[Hyperparameters, training details, compute infrastructure.]

\section{Additional Experiments}

[Extra ablations, sensitivity analyses, qualitative examples.]

\section{Proofs}

[Full proofs of any theoretical results in the main text.]

\end{document}
